%=======================================================================
%  ABBREVIATION CHART  (shared, appended to every chapter)
%  At most 16 rows per tabularx block: tabularx cannot break across a page,
%  so a long block moves whole and dumps a gap at the foot of the previous
%  page. \par\penalty0 between blocks reopens the breakpoint that \vspace
%  would otherwise close.
%=======================================================================
\needspace{6\baselineskip}
\T{Abbreviations used in these notes}

{\color{sub}\footnotesize Only these are ever shortened. Anything not on this list is written out in full.}

\vspace{2pt}
\begin{tabularx}{\textwidth}{@{}L{1.75cm}X@{\hspace{8pt}}L{1.75cm}X@{}}
\toprule
\hrow \thead{Short} & \thead{Means} & \thead{Short} & \thead{Means} \\
\midrule
AI       & artificial intelligence            & ANI      & artificial narrow intelligence \\
AGI      & artificial general intelligence    & ASI      & artificial super intelligence \\
KR       & knowledge representation           & KRR      & knowledge representation and reasoning \\
PEAS     & performance measure, environment, actuators, sensors & GPS & General Problem Solver \\
NLP      & natural language processing        & NLU      & natural language understanding \\
NLG      & natural language generation        & CAPTCHA  & completely automated public Turing test to tell computers and humans apart \\
ML       & machine learning                   & ANN, NN  & artificial neural network \\
MLP      & multilayer perceptron              & BP       & back-propagation \\
SOM      & self-organizing map                & GA       & genetic algorithm \\
FIS      & fuzzy inference system             & ID3      & Iterative Dichotomiser 3 \\
CSP      & constraint satisfaction problem    & DFS      & depth first search \\
BFS      & breadth first search               & DLS      & depth limited search \\
IDS      & iterative deepening search         & UCS      & uniform cost search \\
CNF      & conjunctive normal form            & WFF      & well-formed formula \\
FOL      & first-order logic                  & FOPL     & first-order predicate logic \\
RRS      & resolution refutation system       & BBN      & Bayesian belief network \\
\bottomrule
\end{tabularx}
\par\penalty0
\vspace{4pt}
\begin{tabularx}{\textwidth}{@{}L{1.75cm}X@{\hspace{8pt}}L{1.75cm}X@{}}
\toprule
\hrow \thead{Short} & \thead{Means} & \thead{Short} & \thead{Means} \\
\midrule
ES       & expert system                      & KB       & knowledge base \\
CD       & conceptual dependency              & LISP     & list processing (the language) \\
MOP      & memory organization packet         & isa      & the class-to-class link of a semantic net (\emph{instance} is the individual-to-class one) \\
attr(s)  & attribute(s)                       & repo     & repository \\
freq     & frequency                          & config   & configuration \\
info     & information                        & no.      & number \\
Eg       & example                            & Vs       & compare / differentiate \\
w/, w/o  & with, without                      & $+$Fig   & draw the figure \\
$\rightarrow$ & leads to, therefore           & $=$      & is, equals \\
\bottomrule
\end{tabularx}
\par\penalty0
\vspace{4pt}
\begin{tabularx}{\textwidth}{@{}X@{}}
\toprule
\hrow \thead{Reading the year tags and chips} \\
\midrule
\textbf{Exam-paper month codes}:
Ba $=$ Baishakh $\cdot$ Jth $=$ Jestha $\cdot$ Shr $=$ Shrawan $\cdot$ Bh $=$ Bhadra $\cdot$
Ash $=$ Ashwin $\cdot$ Ka $=$ Kartik $\cdot$ Po $=$ Poush $\cdot$ Ma $=$ Magh $\cdot$
Ch $=$ Chaitra. Years are Bikram Sambat, so \yr{80 Ch} is 2080 Chaitra. \\
\textbf{Bold} year tags are Regular exams, plain are Back exams. A header printing
``Regular / Back'' in one box counts as Regular. \\
\textbf{Three course codes, three faces.} Plain $=$ \textbf{CT 653} (BCT, III-II, 26
papers) $\cdot$ \texttt{monospace} $=$ \texttt{CT 710} (BEI, IV-I, 7 papers) $\cdot$
\textit{italic} $=$ \textit{CT 78506} (BEX, IV-II, 8 papers).
\mk{The same year code names different papers under different codes}: \yr{\textbf{77 Ch}}
and \yr{\textbf{\textit{77 Ch}}} are two different exams, and only the font tells them
apart. \\
\textbf{Tier chips} over all 41 papers, 2069--2082 BS:
\tS{} $=$ topic asked in \textbf{8 or more} papers $\cdot$
\tF{} $=$ \textbf{4 to 7} papers $\cdot$
\tP{} $=$ \textbf{1 to 3} papers but worth $\ge$6 marks. \\
\textbf{Mark chips} such as \m{2+6} are the marks that question actually carried.
\pill{gdL}{gd}{syllabus, no PYQ yet} marks a syllabus topic never yet examined but written
up anyway. \added\ marks anything researched and added that no source note in this folder
covers. \\
\bottomrule
\end{tabularx}
